\documentclass[11pt]{article}

\usepackage{amsmath, amssymb, amsfonts}
\usepackage{geometry}
\geometry{margin=1in}

\title{Supplementary Information: Information Distortion, Corrected Covariance, and Bias}
\author{}
\date{}

\begin{document}

\maketitle

\section*{1. Log-Likelihood and Score Definitions}

Consider a trajectory of $T$ sequential measurement intervals.
The total log-likelihood is defined as

\begin{equation}
\ell(\theta) = \log L(\theta) = \sum_{t=1}^T \ell_t(\theta),
\end{equation}

with total score

\begin{equation}
S(\theta) = \nabla_\theta \ell(\theta)
= \sum_{t=1}^T s_t(\theta),
\qquad
s_t(\theta) = \nabla_\theta \ell_t(\theta).
\end{equation}

Expectations and covariances are taken over independent realizations of full trajectories generated under parameter $\theta$.

\section*{2. Gaussian Fisher Information (GFI)}

At each interval, the predictive likelihood is approximated as Gaussian:

\begin{equation}
y_t \sim \mathcal{N}\!\left(\mu_t(\theta), \sigma_t^2(\theta)\right),
\end{equation}

where $\mu_t(\theta)$ and $\sigma_t^2(\theta)$ are obtained recursively.

For scalar variance $\sigma_t^2$, the per-interval Gaussian Fisher information is

\begin{equation}
\mathbf I_t^{\mathrm{G}}(\theta)
=
\frac{1}{\sigma_t^2(\theta)} \,
\nabla_\theta \mu_t(\theta)
\nabla_\theta \mu_t(\theta)^\top
+
\frac{1}{2\sigma_t^4(\theta)} \,
\nabla_\theta \sigma_t^2(\theta)
\nabla_\theta \sigma_t^2(\theta)^\top.
\end{equation}

The total Gaussian Fisher Information is

\begin{equation}
\mathbf H(\theta)
=
\sum_{t=1}^T \mathbb{E}\!\left[\mathbf I_t^{\mathrm{G}}(\theta)\right]
=
- \mathbb{E}\!\left[\nabla_\theta^2 \ell(\theta)\right].
\end{equation}

This matrix defines the local quadratic geometry implied by the Gaussian predictive model.

\section*{3. Score Covariance Matrix (SCM)}

The Score Covariance Matrix is defined as

\begin{equation}
\mathbf J(\theta)
=
\mathrm{Cov}\!\left(S(\theta)\right).
\end{equation}

This matrix measures the total fluctuation strength of the score process,
including temporal correlations and model inconsistencies.

\section*{4. Information Distortion Matrix (IDM)}

We define the Information Distortion Matrix as

\begin{equation}
\mathbf C(\theta)
=
\mathbf H(\theta)^{-1/2}
\mathbf J(\theta)
\mathbf H(\theta)^{-1/2}.
\end{equation}

$\mathbf C$ is symmetric positive definite and dimensionless.

If the Gaussian predictive model fully captures conditional dynamics,

\begin{equation}
\mathbf J(\theta_0) = \mathbf H(\theta_0),
\qquad
\mathbf C(\theta_0) = \mathbf I.
\end{equation}

\section*{5. Per-Parameter Distortion}

In the original parameter basis, the diagonal elements of $\mathbf C$ quantify
the variance inflation (or deflation) of each parameter relative to the Gaussian Fisher prediction:

\begin{equation}
C_{ii}
=
\frac{\text{real fluctuation variance in direction } i}
     {\text{Gaussian Fisher predicted variance in direction } i}.
\end{equation}

Thus:

\begin{itemize}
\item $C_{ii} = 1$ indicates consistency for parameter $i$.
\item $C_{ii} > 1$ indicates underestimation of uncertainty.
\item $C_{ii} < 1$ indicates overestimation of uncertainty.
\end{itemize}

While $\mathbf C$ contains full directional information, we use its diagonal
as a physically interpretable per-parameter distortion measure.

\section*{6. Distortion-Corrected Covariance (DCC)}

The covariance of the estimator is given asymptotically by

\begin{equation}
\Sigma_{\mathrm{DCC}}
=
\mathbf H^{-1}
\mathbf J
\mathbf H^{-1}
=
\mathbf H^{-1/2}
\mathbf C
\mathbf H^{-1/2}.
\end{equation}

We refer to $\Sigma_{\mathrm{DCC}}$ as the Distortion-Corrected Covariance (DCC).

When $\mathbf C = \mathbf I$, $\Sigma_{\mathrm{DCC}} = \mathbf H^{-1}$.

\section*{7. Distortion-Induced Bias (DIB)}

If the expectation of the total score is nonzero,

\begin{equation}
\mathbf g(\theta) = \mathbb{E}[S(\theta)] \neq 0,
\end{equation}

then the estimator exhibits a first-order bias:

\begin{equation}
\mathbf b_{\mathrm{DIB}}
=
\mathbf H^{-1} \mathbf g(\theta).
\end{equation}

We refer to $\mathbf b_{\mathrm{DIB}}$ as the Distortion-Induced Bias (DIB).

It represents the systematic parameter displacement induced by persistent score drift.

\section*{8. Scaling with Sample Size}

For $N$ independent trajectories:

\begin{equation}
\Sigma_{\mathrm{DCC}} \propto \frac{1}{N},
\end{equation}

while structural bias remains approximately constant.
Thus the relative importance of bias increases with sample size.

\section*{9. Summary}

\begin{align}
\mathbf H &\rightarrow \text{Gaussian Fisher curvature} \\
\mathbf J &\rightarrow \text{Score fluctuation geometry} \\
\mathbf C &\rightarrow \text{Information distortion} \\
\Sigma_{\mathrm{DCC}} &\rightarrow \text{Distortion-Corrected Covariance} \\
\mathbf b_{\mathrm{DIB}} &\rightarrow \text{Distortion-Induced Bias}.
\end{align}

Together, these quantities provide a complete geometric description of
uncertainty and systematic distortion induced by the predictive algorithm.

\end{document}
